% ============================================================
% Kapitel 7: Ausblick
% ============================================================
\chapter{Fazit und Ausblick}

Zusammenfassend ist zu sagen dass On premise lösungen welche lokal auf server betrieben werden in bestimmmten anwendungsfällen eine sehr gut alternative zu cloud basierten lösungen bieten. Robotikanwendungen in innerbetrieblichen Bereich sind oftmals überschaubar genung bezogen auf veränderlichen variabeln, dass die anpassung der hyperparamter eines klein bis mittlgroßen modells überschaubar sind und demnach gut lokal betrieben werden können. Für große Modelle wie die meißsten sie heute kennen sind große datenzentren nötig mit hohen energiekosten welche für einzelen anwendungen nur wenig sinn ergeben. Dennoch ist immer abzugwwägen ob es in Zukunft für ein wachsendes Unternehmen nicht von Vorteil wäre geld in die hand zu nehmen um eine Serverstruktur zu schaffen welche sich vom Cloud geschäft ein stück weit loslöst und eigen LLM auf eigenen Servern betreibt. Zudem wird trotz wachsender Compute möglichkeiten eine linie nicht überschritten werden können laut \ref{} 


\section{Zusammenfassung der Ergebnisse}
% TODO: Kernerkenntnisse der Arbeit

\section{Ausblick auf die Masterarbeit}
Die Masterarbeit knüpft an vielen der hier erarbeiteten Grundlagen direkt an und umfasst weiterfürhende Themen wie vision language action models in der Robotik. Zusammene mit einer weiteren Forschungs und Entwicklungsarbeit zum Thema ROS2 auf einem IGus Rebel 6 DOf wird ein Testaufbau entwickelt welcher konkret Robotik und Large Languange models im Kontext von Open source und on premise evaluiert.
Generell bleit abzuwarten ob der hype um diese neue Technologie den erhofften wandel der gesellschaft standhalten kann oder doch kompromisse eingegangen werden müssen. auf jden fall wurde durch diese Forschungsarbeit ein Grundstein zur Erklärung der Modelle gelegt.
